problem:  
Let \((X,L)\) be a polarized compact Kähler manifold with discrete automorphism group, and \(E\) the Mabuchi functional on the space \(\mathcal{H}\) of smooth Kähler potentials in \(c_1(L)\).  Define for any real parameter \(\lambda>0\) the **entropy‑perturbed Mabuchi functional**
\[
E_{\lambda}(\varphi)=E(\varphi)+\lambda\,\mathrm{Ent}(\omega_{\varphi}^n,\omega_0^n)
=E(\varphi)+\lambda\!\int_X\!\log\!\Big(\frac{\omega_{\varphi}^n}{\omega_0^n}\Big)\omega_{\varphi}^n,
\]
where \(\mathrm{Ent}(\omega_{\varphi}^n,\omega_0^n)\) denotes the relative entropy of volume forms.

Let \(\mathcal{E}^1\) be the finite‑energy completion of \(\mathcal{H}\) with respect to the \(L^1\) Mabuchi metric.  The classical coercivity conjecture asserts that \(E\) is proper (modulo automorphisms) on \(\mathcal E^1\) if and only if \((X,L)\) is uniformly K‑stable.

**Research question.**
1. For fixed \((X,L)\), does there exist a finite constant \(\lambda_0(X,L)\ge0\) such that \(E_\lambda\) is coercive on \(\mathcal E^1\) for every \(\lambda>\lambda_0\)?
2. If such a threshold exists, can \(\lambda_0(X,L)\) be computed (or bounded) in terms of algebraic invariants of \((X,L)\), for example the \(\delta\)‑invariant or the \(\alpha\)‑invariant?
3. Does the behavior of \(\lambda_0(X,L)\) under degenerations (e.g. in test configurations) provide a new quantitative invariant interpolating between analytic coercivity and non‑Archimedean stability?

Equivalently, this asks whether adding a multiple of the entropy functional can *regularize* the Mabuchi energy—forcing convexity or properness beyond the K‑stable range—and whether the minimal such multiple encodes intrinsic algebro‑geometric data.

Why is it a "good" problem:  
This refined formulation eliminates speculative claims while focusing on a precise analytic stability threshold problem.  The question is mathematically sound and explores a deep interface between scalar curvature geometry, information‑theoretic convexity, and algebraic stability theory.  It asks whether the entropy term can act as a universal analytic stabilization mechanism for the Mabuchi functional, which would reveal a tunable “regularization scale’’ linking nonlinear geometric analysis and K‑stability.  Determining the existence or nature of the threshold \(\lambda_0(X,L)\) would open a new axis of quantitative comparison between coercivity in pluripotential theory and δ‑invariants in algebraic geometry, potentially leading to novel criteria for canonical metrics.  The statement is simple yet demands profound understanding of convexity in infinite‑dimensional spaces and algebro‑geometric measures of positivity—therefore it is a genuinely **good** research‑level problem.